\section{Experiments}

\subsection{Experimental Setup}


\begin{figure*}[t]
\centering
\includegraphics[width=1.0\textwidth]{./assets/demo_main.pdf}
\vspace{-15pt}
\caption{\textbf{Qualitative results of our model.} 
The first row shows the 4D point cloud generated by our 4D-STraG. The second and third rows show the videos rendered by our 4D-ViSM under two distinct, user-defined camera trajectories. }
\vspace{-5pt}
\label{fig:demo_main}
\end{figure*}

\noindent\textbf{Implementation Details.} 
Trained on TrajScene-60K dataset, our model generates  $512\times368$ videos with a length of 49 frames. For the 4D-STraG module, we performed full-parameter training based on Wan2.1-14B~\cite{wan2025wan,wan_fun}. The trajectory encoder and decoder are both shallow ResNets. We first trained its tracking components and VAE decoder for 5k steps, then its DiT for 2k steps using OmniMAE features for motion conditioning. The 4D-ViSM module finetuned Wan2.1-14B~\cite{wan2025wan,wan_fun_inp} with LoRA for 10k steps. We used AdamW~\cite{adamw} with a learning rate of $2\times 10^{-5}$. All experiments ran on four NVIDIA H20 GPUs.



\subsection{Qualitative Results}
\label{sec:qualitative}

For qualitative validation of our model effectiveness, we conduct comparative analyses in Figure~\ref{fig:demo_main} and Figure~\ref{fig:compare}. Figure~\ref{fig:demo_main} presents two representative scenes, with the first row displaying the point cloud outputs of 4D-STraG. It can be observed that the 4D point clouds maintain structural consistency and exhibit reasonable motion over time, with high detail completeness. For each scene, we define two camera trajectories and perform rendering using 4D-ViSM. The results demonstrate that the rendered images not only preserve the geometric accuracy of the original structure but also effectively align with the camera trajectories, reflecting strong visual consistency. Figure~\ref{fig:compare} further compares the visual outcomes of several existing methods, including 4Real, DimensionX, Gen3C, and Free4D. In comparison to these approaches, our method generates more diverse and realistic motion, validating the advantage of the our strategy in ensuring both structural rationality and motion plausibility for 4D generation. 


\begin{figure*}[t!]
\centering
\includegraphics[width=1.0\textwidth]{./assets/compare2.pdf}
\vspace{-16pt}
\caption{\textbf{Qualitative comparison with baseline methods.} For each sample, the first row shows the baseline results while the second row presents our \ours results. The first column displays the input image and text prompt.}
\label{fig:compare}
\vspace{-4pt}
\end{figure*}


\begin{table*}[t!]
\footnotesize  
\centering
\caption{\textbf{Quantitative comparison on VBench.} Higher values are better. The best results in each comparison group are marked in \textbf{bold}. Paradigm: Opt. = Optimization-based, Learn. = Learning-based, TF = Training-free.}
\vspace{-3mm}
\label{tab:vbench_results}
\renewrobustcmd{\bfseries}{\fontseries{b}\selectfont} 
\renewrobustcmd{\mdseries}{\fontseries{m}\selectfont}
\setlength{\tabcolsep}{6pt}
\small
\resizebox{\textwidth}{!}{%
\begin{tabular}{c|l|l|cccccc}
\toprule
\textbf{Exp.} & \multirow{2}{*}{\textbf{Model / Metrics}}
& \multirow{2}{*}{\textbf{Paradigm}}
& \textbf{Subject} & \textbf{Background}
& \textbf{Motion} & \textbf{Dynamic}
& \textbf{Aesthetic} & \textbf{Imaging} \\
 \textbf{No.} & & & {\textbf{Consistency}} & {\textbf{Consistency}}
& \textbf{Smoothness} & \textbf{Degree}
& \textbf{Quality} & \textbf{Quality} \\
\midrule
% --- First Comparison Group ---
\multirow{2}{*}{\textbf{I}} & 4Real~\cite{4real} & Opt.   &  \bfseries  0.9329 & \bfseries  0.9709 & 0.9664 & 0.7708 & 0.4938 & 0.5095  \\
 & \ours(Ours) & Learn.    & 0.8752 &  0.9364      &    \bfseries 0.9682      &  \bfseries  1.0000       &   \bfseries  0.5613      &    \bfseries   0.6230    \\
\midrule
% --- Second Comparison Group ---
\multirow{3}{*}{\textbf{II}} & GenXD~\cite{genXD} & Learn.   & 0.8042 & 0.8789 & 0.9030 & \bfseries 1.0000 & 0.4077 & 0.5209 \\
 & DimensionX~\cite{dimensionX} & Learn. &  0.7553 &  0.8481 & \bfseries 0.9827 & \bfseries 1.0000 &  0.4634 &  0.5545 \\
 & \ours(Ours) & Learn.    & \bfseries 0.8241 & \bfseries 0.9044 &  0.9760 & 0.9500 & \bfseries 0.4820 & \bfseries 0.5828  \\
\midrule
% --- Third Comparison Group ---
\multirow{3}{*}{\textbf{III}} & Free4D~\cite{free4d} & TF     & 0.7899 & 0.8883 & 0.9797 & \bfseries 1.0000 & 0.3607 & 0.3562  \\
 & Gen3C~\cite{gen3c} & Learn.      & 0.8112 & 0.8871 & \bfseries  0.9845 & 0.9940 & 0.3812 & 0.4814  \\
 & \ours(Ours) & Learn.       & \bfseries 0.8339 & \bfseries 0.9065 & 0.9773 & 0.9000 & \bfseries 0.4820 & \bfseries 0.5939  \\
\bottomrule

\end{tabular}
}
\vspace{-4mm}
\end{table*}


\subsection{Quantitative Results}
\label{sec:quantitative}
\noindent\textbf{Video Quality Benchmarking.} Following prior 4D generation works~\cite{free4d}, we use VBench~\cite{vbench} for a fair and extensive video-quality comparison against a wider range of baselines. A more detailed 4D geometric analysis is provided in the supplementary material. We compare it with leading 4D generation methods: 4Real, GenXD, DimensionX, Gen3C, and Free4D. Following the evaluation protocol in Free4D~\cite{free4d}, we structured our comparisons into three groups based on model availability and technical constraints, with each group evaluated under appropriate camera trajectories. In Group I, we compared against the closed-source 4Real using simple trajectories, as we could only access their official demonstration videos for evaluation. Group II benchmarks single-image 3D models (GenXD, DimensionX's S-Director) with moderate 90° left rotations, due to limited trajectory support in their open-source code. For Group III, we evaluated against other 4D generation models using complex trajectories, including upward, forward, leftward, rightward, and downward 90° movements, to assess performance under challenging camera motions. For Groups II and III, we used 200 held-out samples from WebVid-10M that do not overlap with our TrajScene-60k training subset to ensure fair evaluation in diverse scenarios.  As shown in Table \ref{tab:vbench_results}, \ours consistently achieved promising results across groups. Notably, it demonstrated superior Dynamics, Aesthetic, and Imaging Quality compared to 4Real, outperformed 3D reconstruction baselines in consistency and visual quality, and maintained a pronounced lead in Aesthetic and Imaging Quality against 4D generation models under complex trajectories. Furthermore, consistent improvements across all groups demonstrate our robust generalization.



\begin{table*}[t!]\small
\small  
\centering
\caption{\textbf{Quantitative evaluation on 4D Generation Consistency via VLM-based assessment.} Higher values are better.}
\vspace{-3mm}
\label{tab:4d_video_results}
\sisetup{detect-weight, mode=text} % Setup siunitx to detect bold weight
\renewrobustcmd{\bfseries}{\fontseries{b}\selectfont} % Make bfseries robust for siunitx
\renewrobustcmd{\mdseries}{\fontseries{m}\selectfont}
\setlength{\tabcolsep}{3pt} %
\resizebox{\textwidth}{!}{%
\begin{tabular}{c|l|ccccc|c}
\toprule
\textbf{Exp.} & \multirow{2}{*}{\textbf{Model / Metrics}} 
& \textbf{3D Geometric} & \textbf{Temporal Texture} 
& \textbf{Subject Identity} & \textbf{Motion Geometry} 
& \textbf{Background Stability} & \textbf{Average} \\
 \textbf{No.} & & {\textbf{Consistency}} & {\textbf{Stability}}  
& {\textbf{Preservation}} & {\textbf{Coupling}} 
& {\textbf{Consistency}} & \textbf{Score} \\
\midrule
% --- First Comparison Group ---
\multirow{2}{*}{\textbf{I}} & 4Real~\cite{4real}   & 3.04 & 3.49 & 4.11 & 2.64 & 3.49 & 3.35 \\
 & \ours (Ours)    & \bfseries 3.46 & \bfseries 3.88 & \bfseries 4.42 & \bfseries 3.25 & \bfseries 3.96 & \bfseries 3.80 \\
\midrule
% --- Second Comparison Group ---
\multirow{3}{*}{\textbf{II}} & GenXD~\cite{genXD}   & 2.43 & 2.66 & 3.09 & 2.38 & 2.65 & 2.64 \\
 & DimensionX~\cite{dimensionX} & 2.58 & 2.74 & 3.35 & 2.53 & 2.79 & 2.80 \\
 & \ours (Ours)    & \bfseries 3.53 & \bfseries 3.92 & \bfseries 4.33 & \bfseries 3.45 & \bfseries 3.93 & \bfseries 3.83 \\
\midrule
% --- Third Comparison Group ---
\multirow{3}{*}{\textbf{III}} & Free4D~\cite{free4d}     & 1.13 & 1.30 & 1.37 & 1.17 & 1.17 & 1.23 \\
 & Gen3C~\cite{gen3c}      & 1.99 & 2.26 & 2.54 & 2.01 & 2.31 & 2.22 \\
 & \ours (Ours)       & \bfseries 3.47 & \bfseries 3.88 & \bfseries 4.25 & \bfseries 3.35 & \bfseries 3.85 & \bfseries 3.76 \\
\bottomrule
\end{tabular}
}
\vspace{-4mm}
\end{table*}

\noindent\textbf{4D Consistency Analysis.} Because standard metrics like VBench insufficiently penalize 3D-specific artifacts (e.g., geometric drift or texture swimming), we introduce a VLM-based protocol for comprehensive 4D evaluation. For videos generated across the aforementioned three groups, we uniformly sample 8 frames and prompt Qwen2.5-VL-72B-Instruct~\cite{bai2025qwen2} to rate them on a 1-5 scale across five criteria: 3D Geometric Consistency, Temporal Texture Stability, Subject Identity, Motion-Geometry Coupling, and Background Stability. We provide the exact scoring prompt and frame sampling strategy in the supplementary material for full reproducibility. As Table \ref{tab:4d_video_results} shows, \ours consistently outperforms all baselines. Notably, our significant margins in Motion-Geometry Coupling and Temporal Texture Stability validate \ours's superior spatial-temporal coherence and structural preservation during complex camera motions.

Moreover, to explicitly quantify structural stability, we conduct two rigorous geometry-aware evaluations: (1) \textbf{Avg. Trajectory Error}, measuring the $L_2$ distance between point trajectories tracked by DELTA in generated videos versus their ground-truth counterparts (evaluated on all generated samples from Groups II and III following DELTA's default tracking settings); and (2) \textbf{Avg. 3D Reprojection Error}, which follows WorldScore~\cite{zhang2025world} by utilizing DROID-SLAM~\cite{teed2021droid} to gauge the geometric deviation across all co-visible pixels in consecutive frames. As shown in Table~\ref{tab:3d_sub}, \ours achieves competitive 3D consistency. Notably, our method reduces the trajectory error by a substantial margin while maintaining highly competitive or superior reprojection errors. These results explicitly validate \ours's exceptional pixel-level tracking accuracy and robust rigid 3D spatial coherence.

\noindent\textbf{Inference Efficiency.} As Table~\ref{tab:runtime_comparison_sub} shows, unlike many optimization-based baselines that require hours of computation, \ours achieves highly competitive inference speed. On a single A100 GPU, it synthesizes 49-frame, $512\times368$ sequences in just 6 minutes. While GenXD is faster (2 mins), it yields significantly shorter, lower-resolution outputs. Crucially, we are $5\times$ faster than Free4D under equivalent resolutions. This remarkable efficiency allows our pipeline to inherently support longer-duration video synthesis.


\begin{table}[t]
\centering
\caption{\textbf{Runtime and error comparison.} Left: runtime on a single NVIDIA A100 (avg.\ 100 samples). Right: trajectory \& reprojection errors (lower is better)}
\label{tab:runtime_and_errors}
\setlength{\tabcolsep}{4pt}
\renewcommand{\arraystretch}{1.06}
  \vspace{-14pt} % 强制顶对齐
\begin{subtable}[t]{0.45\textwidth}
  \centering
  \caption{Mean Trajectory \& Reprojection Errors}
  \label{tab:errors_sub}
  \small
  \vspace{-10pt} % 强制顶对齐
  % 使用 siunitx 的 S 列可以在有具体数值时改进对齐；此处用普通列占位
  \resizebox{\linewidth}{!}{%
  \begin{tabular}{@{}c |l| c c@{}}
    \toprule
    \label{tab:3d_sub}
    \textbf{Exp.} & \textbf{Model} & \textbf{Trajectory Error} $\downarrow$ & \textbf{Reprojection Error} $\downarrow$ \\
    \midrule
    \multirow{4}{*}{\textbf{i}} 
      & \textcolor{gray}{Ground truth}           & \textcolor{gray}{0.000} & \textcolor{gray}{0.301} \\
      & GenXD~\cite{genXD}           & 0.236 & 1.891 \\
      & DimensionX~\cite{dimensionX} & 0.465 & \textbf{0.602} \\
      & \ours (Ours)                 & \textbf{0.058} & 0.614 \\
    \midrule
    \multirow{4}{*}{\textbf{ii}} 
      & \textcolor{gray}{Ground truth}           & \textcolor{gray}{0.000} & \textcolor{gray}{0.284} \\
      & Free4D~\cite{free4d}         & 0.252 & 0.652 \\
      & Gen3C~\cite{gen3c}           & 0.197 & 0.755 \\
      & \ours (Ours)                 & \textbf{0.042} & \textbf{0.639} \\
    \bottomrule
  \end{tabular}
  }% end resizebox
\end{subtable}
\hfill
\begin{subtable}[t]{0.48\textwidth}
  \centering
  \caption{Runtime Analysis(on a single A100)}
  \label{tab:runtime_comparison_sub}
  \small
  \vspace{-10pt} % 强制顶对齐
  \resizebox{\linewidth}{!}{%
  \begin{tabular}{@{}l| c c c c@{}}
    \toprule
    \textbf{Method} & \textbf{Year} & \textbf{Resolution} & \textbf{Frames} & \textbf{Time} \\
    \midrule
    4Dfy~\cite{4dfy}     & CVPR'24    & $256\times256$  & $-$   & 10 h   \\
    4Real~\cite{4real}   & NeurIPS'24 & $256\times144$  & 8     & 1.5 h  \\
    Gen3C~\cite{gen3c}   & CVPR'25    & $1280\times704$ & 121   & 50 min \\
    GenXD~\cite{genXD}   & ICLR'25    & $256\times256$  & 12    & 2 min  \\
    Free4D~\cite{free4d} & ICCV'25    & $512\times368$  & 16    & 30 min \\
    \midrule
    \ours (Ours)        & $-$        & $512\times368$  & 49    & 6 min  \\
    \bottomrule
  \end{tabular}
  }% end resizebox
\end{subtable}%

\vspace{-2pt}
\end{table}

\begin{table}[t!]
\centering
\caption{\textbf{Component-wise Ablation Study.} 
We analyze the impact of depth normalization, latent design, MPM structure, and training data variations. 
Best results per metric are highlighted in \textbf{bold}.}
\vspace{-4pt} % 强制顶对齐
\label{tab:ablation_transposed}
\footnotesize
\setlength{\tabcolsep}{4pt}
\renewcommand{\arraystretch}{1.05}
\resizebox{\textwidth}{!}{%
\begin{tabular}{l|c c c c c c c c}
\toprule
\multirow{2.5}{*}{\textbf{Metric}} &
\multirow{2.5}{*}{\makecell{\textbf{w/o Depth}\\\textbf{Norm}}} &
\multirow{2.5}{*}{\makecell{\textbf{w/o Depth}\\\textbf{Latents}}} &
\multicolumn{2}{c}{\textbf{MPM Module}} &
\multicolumn{3}{c}{\textbf{Training Data}} &
\multirow{2.5}{*}{\textbf{\ours}} \\
\cmidrule(lr){4-5} \cmidrule(lr){6-8}
& & & \textbf{w/o MPM} & \textbf{w/o Patch Feat.} 
& \textbf{Small Scale} & \textbf{Low Quality} & \textbf{Noise} & \\
\midrule
\textbf{Consistency↑} 
& 0.8604 & 0.8567 & 0.8650 & \textbf{0.8743} 
& 0.8567 & 0.8549 & 0.8685 & 0.8702 \\

\textbf{Dynamic↑}     
& 0.8850 & 0.8500 & 0.8500 & 0.8840 
& 0.8920 & 0.8850 & 0.8967 & \textbf{0.9000} \\

\textbf{Aesthetic↑}   
& 0.4672 & 0.4738 & 0.4806 & 0.4754 
& 0.4791 & 0.4654 & 0.4771 & \textbf{0.4820} \\
\bottomrule
\end{tabular}
}
\vspace{-4pt} % 强制顶对齐
\end{table}





\begin{figure*}[t]
\centering
\includegraphics[width=1.0\textwidth]{./assets/ablation_all.pdf}
\vspace{-10pt}
\caption{\textbf{Ablation studies on normalization methods and module components.} (Rows 1-2) Depth-guided motion normalization stabilizes 4D point cloud generation. (Rows 3-6) Removing the MPM module reduces motion magnitude while excluding depth guidance breaks structural motion consistency, validating our design choices.}
\vspace{-3pt}
\label{fig:ablation_module}
\end{figure*}


\begin{figure}[t]
  \centering
  % \vspace{-10mm}
  \includegraphics[width=\linewidth]{assets/failure.pdf}
  \vspace{-8pt}
  \caption{Limitations: Extreme lighting and unstructured crowds.}
   \label{fig:failure}
   \vspace{-12pt}
\end{figure}


\subsection{Ablation Studies}

Table~\ref{tab:ablation_transposed} presents a comprehensive ablation study on VBench under the Group III setting. Removing depth-guided normalization drops Consistency to 0.8604 and Aesthetic quality to 0.4672, while excluding depth latents degrades both Consistency and Dynamic scores, as the model struggles to maintain motion coherence. Removing MPM reduces Dynamic from 0.9 to 0.85; replacing its patch-level features with a global \texttt{[CLS]} token slightly improves Consistency but degrades Dynamic and Aesthetic quality, confirming that fine-grained patch-level conditioning is essential for high-fidelity 4D generation. For training data, reducing scale to a 1k subset or using unfiltered 60k samples both cause significant performance drops, validating that large-scale, high-quality data is foundational. Notably, adding random noise during training yields only negligible degradation, demonstrating robustness to noisy pseudo-GT supervision. Our full \ours achieves optimal performance across all metrics, validating the synergistic contribution of each component.

Row 1-6 in Figure~\ref{fig:ablation_module} visually demonstrate these findings. Baseline min-max normalization (rows 1-2) yields unstable point clouds with exaggerated movements, particularly in scenes with large depth variations. Our depth-guided approach effectively resolves this issue. Further visualizations (rows 3-6) show that without MPM or depth guidance, the model fails to preserve object structure or maintain motion consistency across different parts, confirming the critical role of depth in ensuring structurally coherent point trajectories.

\subsection{Discussions}




\noindent\textbf{Geometry-Conditioned vs.\ Sequential 4D Generation.} To validate our geometry-conditioned trajectory generation paradigm, we compare against two sequential baselines: Wan2.1-I2V~\cite{wan2025wan} followed by either DELTA~\cite{delta} tracking or VGGT~\cite{vggt} reconstruction. As shown in Figure~\ref{fig:vggt_comp}, sequential pipelines suffer from severe error accumulation: since video diffusion models lack explicit 3D constraints, frame-level appearance inconsistencies are misinterpreted as 3D motion by downstream modules, producing spurious background drift and geometric fragmentation. In contrast, our approach conditions trajectory generation directly on the initial point cloud geometry, enabling the diffusion model to produce temporally coherent, geometry-consistent 4D trajectories without relying on error-prone intermediate video synthesis.


\noindent\textbf{Limitations.}
As shown in Fig.~\ref{fig:failure}, \ours struggles with extreme photometric conditions (e.g., harsh backlighting) and highly unstructured dynamics (e.g., chaotic crowds). In these complex edge cases, off-the-shelf monocular depth estimators often degrade, introducing noise into the initial geometric prior. Since our dense trajectory generation is conditionally bounded by the quality of this initial geometry, the resulting motion can become less reliable. Integrating more robust visual foundation models to handle such extreme physical conditions remains an important direction for future work.



\begin{figure*}[t]
\centering
\includegraphics[width=1\textwidth]{./assets/vggt_comp.pdf}
\vspace{-10pt}
\caption{\textbf{Geometry-conditioned 4D-STraG (top row) vs.\ sequential pipelines.} Rows 2--4 show cascaded baselines: Wan2.1-I2V video generation followed by DELTA tracking or VGGT reconstruction. By conditioning trajectory generation directly on initial point cloud geometry, our method achieves superior spatio-temporal coherence, whereas sequential approaches suffer fragmentation and drift from error accumulation. All results are rendered from a fixed camera viewpoint for fair comparison.}
\vspace{-4pt}

\label{fig:vggt_comp}
\end{figure*}